% Generated from lab3-report.md; SHA256=045f0e83970d88b31a9b59e8387dad32bb3aa22b6adf0a91339df1eb66f86fea
% -*- coding: utf-8 -*-
% 基于所提供的南开大学实验报告模板；编译器：XeLaTeX。
% 保留提供模板的版心、封面字体、校徽、水印和页眉页脚；正文由Markdown生成。
\documentclass[UTF8,a4paper,10pt]{ctexart}
\usepackage[left=3.17cm,right=3.17cm,top=2.74cm,bottom=2.74cm]{geometry}
\usepackage{amsmath,amssymb,graphicx,float,booktabs,longtable,array,tabularx}
\usepackage{caption,multirow,xcolor,listings,enumitem,needspace}
\usepackage{fancyhdr,tikz,eso-pic}
\usepackage[colorlinks,linkcolor=black,urlcolor=blue,citecolor=black]{hyperref}
% 有 Noto 字体时使用模板字体，否则使用 ctex 默认字体。
\IfFontExistsTF{Noto Serif CJK SC}{\setCJKmainfont{Noto Serif CJK SC}}{}
\newcommand{\yihao}{\fontsize{26pt}{36pt}\selectfont}
\newcommand{\erhao}{\fontsize{22pt}{28pt}\selectfont}
\newcommand{\xiaoer}{\fontsize{18pt}{18pt}\selectfont}
\newcommand{\sanhao}{\fontsize{16pt}{24pt}\selectfont}
\newcommand{\xiaosan}{\fontsize{15pt}{22pt}\selectfont}
\newcommand{\sihao}{\fontsize{14pt}{21pt}\selectfont}
\newcommand{\banxiaosi}{\fontsize{13pt}{19.5pt}\selectfont}
\newcommand{\xiaosi}{\fontsize{12pt}{18pt}\selectfont}
\newcommand{\dawuhao}{\fontsize{11pt}{11pt}\selectfont}
\newcommand{\wuhao}{\fontsize{10.5pt}{15.75pt}\selectfont}
\newcommand{\tabincell}[2]{\begin{tabular}{@{}#1@{}}#2\end{tabular}}
\ctexset{
 section={name={,、},number=\chinese{section}},
 subsection={name={（,）},number=\chinese{subsection}},
 subsubsection={name={,.},number=\arabic{subsubsection}}
}
\pagestyle{fancy}
\fancyhf{}
\fancyhead[L]{\kaishu\small Lab3：中断与异常处理}
\fancyhead[R]{\kaishu\small 操作系统实验报告}
\fancyfoot[C]{\thepage}
\setlength{\headheight}{15pt}
\renewcommand{\headrulewidth}{0.1pt}
\renewcommand{\footrulewidth}{0pt}
\newcommand{\HRule}{\rule{\linewidth}{0.5mm}}
\newcommand{\HRulegrossa}{\rule{\linewidth}{1.2mm}}
\definecolor{nkupurple}{RGB}{112,35,84}
\definecolor{codebg}{RGB}{247,248,250}
\lstset{
 escapeinside={(*@}{@*)},
 breaklines=true,basicstyle=\ttfamily\footnotesize,
 keywordstyle=\color{blue!70}\bfseries,
 commentstyle=\color{gray},stringstyle=\ttfamily,
 columns=fullflexible,keepspaces=true,showstringspaces=false,
 numbers=left,numberstyle=\tiny\color{blue!50},numbersep=6pt,
 frame=trbl,rulecolor=\color{gray!40},backgroundcolor=\color{codebg},
 xleftmargin=1.5em,tabsize=4,aboveskip=0.8em,belowskip=0.8em
}
\setlength{\parindent}{2em}
\setlength{\parskip}{0.25em}
\setlength{\emergencystretch}{3em}
\setlist{nosep,leftmargin=2em}
% 透明紫色水印：12% 不透明度，置于前景以透过代码背景显示。
\newcommand{\watermark}{\AddToShipoutPictureFG{%
 \AtPageLowerLeft{\includegraphics[width=\paperwidth,height=\paperheight]{NKU-watermark.pdf}}}}
\hypersetup{pdftitle={Lab3：中断与异常处理}}

\usepackage{xurl}
\urlstyle{tt}

\IfFontExistsTF{KaiTi}{\setCJKfamilyfont{zhkai}[AutoFakeBold=2]{KaiTi}}{}

\begin{document}
\begin{titlepage}
\centering
\includegraphics[width=0.8\textwidth]{NKU.png}\\[0.8cm]
{\Huge\kaishu\bfseries 南\quad 开\quad 大\quad 学\par}
\vspace{0.9cm}
{\huge\kaishu\bfseries 计\ 算\ 机\ 学\ 院\par}
\vspace{0.7cm}
{\Large\bfseries 操作系统第三次实验报告\par}
\vspace{0.8cm}
\HRule\par\vspace{0.7cm}
{\LARGE\bfseries 中断与异常处理\\[0.3cm]Trap 现场恢复及定时中断\par}
\vspace{0.6cm}\HRule\par\vspace{1.2cm}
{\large\kaishu
\begin{tabular}{rl}
组长：&李珉熠\quad 2413982\\[0.5cm]
组员：&吴宇石\quad 2411080\\[0.5cm]
组员：&赵亚鑫\quad 2413715
\end{tabular}\par}
\vfill
{\Large 2026年11月2日\par}
\end{titlepage}
\pagenumbering{Roman}
\renewcommand{\abstractname}{\kaishu\sihao\textbf{摘要}}
\begin{abstract}
我们在Lab2基础上完成中断与异常处理机制，包括Trap入口、现场恢复、定时事件重新设置，以及非法指令和断点处理。课程三项评分分别为35、25和40分，实测总分100/100。报告说明实现约束、提示词、真实编译错误和验证范围。
\par\noindent\textbf{关键词：}Trap；上下文；定时中断；异常；RISC-V
\end{abstract}
\clearpage
\tableofcontents
\clearpage
\pagenumbering{arabic}
\watermark


\section{基本信息与分工}

实验日期：2026-11-02。我们基于课程lab3.zip完成实现，使用Codex辅助分析、实现、排错和文档整理，以实际编译、QEMU运行及GDB观察核对结果。以下为小组模块分工。


{\small
\setlength{\tabcolsep}{3pt}
\renewcommand{\arraystretch}{1.22}
\begin{longtable}{>{\raggedright\arraybackslash}p{0.25\linewidth}>{\raggedright\arraybackslash}p{0.68\linewidth}}
\toprule
成员 & 模块 \\
\midrule\endfirsthead
成员 & 模块 \\
\midrule\endhead
李珉熠 2413982（组长） & 环境、Lab2移植、构建与报告汇总 \\
吴宇石 2411080 & Trap入口、汇编恢复和布局分析 \\
赵亚鑫 2413715 & 中断异常处理、评分与调试分析 \\
\bottomrule
\end{longtable}
}

\section{目的、环境与整体逻辑}

我们使用Ubuntu虚拟机、RISC-V GCC 15.1.0、QEMU 8.2.2和OpenSBI 1.3，实验目录为/home/lmy/os/lab3。


Trap是中断与异常的统称。中断是与当前指令执行异步的事件，例如定时器到期；异常由当前指令执行引发，例如非法指令或断点。上下文是被打断时的寄存器状态，trapframe是保存在栈上的现场结构。


整体顺序如下：


\begin{lstlisting}
初始化DTB和物理内存 → 设置stvec统一入口
→ 设置时钟事件和中断使能 → 发生Trap
→ __alltraps / SAVE_ALL → trap(tf) → trap_dispatch
→ interrupt_handler或exception_handler
→ __trapret / RESTORE_ALL → sret返回原执行流
\end{lstlisting}

stvec保存S态Trap入口地址；S态为监管者模式，即本次内核运行的权限级别。硬件设置sepc、scause及相关状态，通用寄存器则由入口汇编保存，不能认为硬件自动完成全部现场保存。


\section{Lab2基础与启动兼容}

我们移植page\_init和Best-Fit，保留Lab3中alloc\_pages、free\_pages和nr\_free\_pages的局部中断保护包装。分配前记录原中断状态并关闭本核中断，操作后恢复原状态，不能无条件重新开中断。这避免本核中断处理打断链表修改，但不提供多核互斥。


物理内存通过DTB取得，描述符与保留区沿用Lab2设计。我们补充PHYSICAL\_MEMORY\_END=0x88000000、KERNEL\_BEGIN\_PADDR=0x80200000两个移植需要的常量。QEMU加载改用-kernel，启动根表增加低地址过渡映射；工具链明确-march=rv64gc、-mabi=lp64d，并使用LF换行。


\section{模块一：Trap入口和现场恢复}

我们补全trapentry.S中的RESTORE\_ALL及\_\_alltraps调用部分，保留SAVE\_ALL布局、trapframe结构和trap函数签名。


RV64下每个槽位8字节，32个通用寄存器槽位加4个控制状态槽位，共36×8=288字节。控制状态依次是sstatus、sepc、stval的保存值和scause。CSR是控制与状态寄存器；stval提供异常附加信息，scause记录原因，sepc记录返回执行位置。


SAVE\_ALL先把原sp放到sscratch，再将sp下移288字节，原sp最终保存在第2号槽位。sp是栈指针，sscratch是辅助暂存CSR。因此我们先读出并恢复sstatus、sepc，再恢复通用寄存器，最后恢复sp，保证所有现场读取仍以当前frame为基址。cause和badvaddr供处理与诊断，不必像sepc一样写回才能返回。


\_\_alltraps用mv a0,sp将frame指针作为第一个参数，call trap调用C函数，函数返回后进入\_\_trapret。a0是RISC-V调用约定中的首参数寄存器，ra保存函数返回地址；旧a0、ra已在frame中，最终一并恢复。


sret是S态Trap返回指令，根据SPP恢复特权级、根据SPIE恢复中断使能，并将sepc作为继续执行的位置。SPP为之前的权限级别，SPIE为之前的中断使能状态。本实验主要由S态返回S态，不是每次都返回用户态。


\section{模块二：中断与异常处理}

\subsection{时钟中断}

trap\_dispatch通过cause最高位区分中断和异常，我们在interrupt\_handler中清除该位后按原因分类。S态时钟原因编号为5。


每次时钟中断调用clock\_set\_next\_event，ticks加1；当ticks是TICK\_NUM=100的整数倍时调用既有print\_ticks。clock\_set\_next\_event通过SBI接口设置未来事件；SBI是内核请求M态固件服务的接口。timebase=100000配合10MHz时间计数器，约每10ms设置下一次事件，实际间隔受处理开销影响，不能等同精确CPU周期。


指导书引用SHUTDOWN\_AFTER，但下载框架未给出定义或明确数值。我们在允许实现区域定义可覆盖默认值10，含义是打印10次100 ticks后调用sbi\_shutdown，而非声称老师规定必须10次。正常运行预期1000次中断后关机；时钟评分在200次时已验证持续返回。


其他软件、定时和外部中断按类型输出提示；未知原因打印编号及frame。外部中断来自设备，软件中断用于软件请求，本实验不实现完整设备中断应答协议。


\subsection{非法指令与断点异常}

我们识别CAUSE\_ILLEGAL\_INSTRUCTION和CAUSE\_BREAKPOINT，输出Illegal instruction at或Breakpoint at。读取epc指向指令的低16位：低2位不全为1，按2字节跳过；否则低5位不全为1，按4字节跳过；更长编码在本教学实现中停止并报告不支持，不猜测跳过。


这适用于本实验RV64GC的16/32位指令。若不修改epc，返回后会再次执行同一异常指令，形成循环。异步时钟中断不主动增加epc，以免跳过正常代码。其他异常按框架要求打印frame，不宣称已经恢复缺页等故障；当前读取故障指令的方式也不适用于任意不可访问epc。


\section{Spec与两份Prompt}

我们将本次协作需求整理为下面两份提示，AI作为辅助，生成结果通过框架约束和实测核对。


\subsection{模块一Spec与Prompt}

Spec：输入现有SAVE\_ALL和36槽位frame；输出与其严格对应的恢复汇编，以及将frame传入trap的调用逻辑；不得修改保存布局、结构体或函数签名。


Prompt：阅读trap.h、trapentry.S和RISC-V调用约定，仅补全RESTORE\_ALL和\_\_alltraps留空处。逐项列出32个通用寄存器及4个状态槽位的偏移，识别原sp经sscratch保存到第2槽位的特殊关系。先恢复返回需要的CSR和通用寄存器，最后恢复sp，避免破坏frame访问基址。以a0传入当前frame，调用trap后接入既有\_\_trapret和sret；保持16字节栈对齐。解释每步为何与SAVE\_ALL对应，并保留lab3注释、替换学号标记。


\subsection{模块二Spec与Prompt}

Spec：接收frame，根据原因分类；定时器持续重设和计数；异常按实际16/32位长度推进epc；保留其余框架函数和课程评分条件。


Prompt：在trap.c规定区域实现interrupt\_handler和exception\_handler。时钟事件每次重设，ticks每次加1，每100次使用已有函数输出，定义并说明可配置关机次数默认值，保证评分可到200次。其他中断分类输出，未知原因打印frame。非法指令和断点按要求输出地址，通过指令低位编码判断2/4字节长度后推进epc，未知异常仅打印现场，异步中断不增加epc。使用现有sbi\_shutdown；不要修改trap\_dispatch、trap签名、SAVE\_ALL或课程测试。


\section{实际排错与验证}

我们编译时实际遇到两处错误：当前汇编器拒绝旧CSR名sbadaddr，因此在汇编文件头用预处理别名映射到同一CSR的现代名称stval，保持SAVE\_ALL主体不变；移植page\_init后缺少两个内存布局常量，核对Lab2与Lab3映射后补充。启动加载参数和过渡映射问题则是根据前次实验经验预先适配，没有将其写成本次已经出现的运行故障。


课程测试源码和注入脚本保留原样，评分脚本只适配QEMU加载方式，未改变通过条件及分数。


{\small
\setlength{\tabcolsep}{3pt}
\renewcommand{\arraystretch}{1.22}
\begin{longtable}{>{\raggedright\arraybackslash}p{0.25\linewidth}>{\raggedright\arraybackslash}p{0.29\linewidth}>{\raggedright\arraybackslash}p{0.37\linewidth}}
\toprule
验证 & 结果 & 能支持的结论 \\
\midrule\endfirsthead
验证 & 结果 & 能支持的结论 \\
\midrule\endhead
现场恢复及断点返回 & 35/35 & 课程预置的多组寄存器经ebreak后保持正确 \\
非法指令及兜底处理 & 25/25 & 4字节非法指令可返回，软件中断与未知原因输出符合要求 \\
持续时钟 & 40/40 & 初始化计数、重新设定和反复返回支持运行到200次 \\
总分 & 100/100 & 本次课程三项检查均通过 \\
普通启动 & 10次100 ticks后正常退出 & 非评分构建的时钟与关机路径可运行 \\
客体GDB & frame=288字节，cause=0x8000000000000005 & 首次时钟进入trap时ticks=0，之后进入\_\_trapret \\
\bottomrule
\end{longtable}
}

完整的16位非法指令、全部扩展指令、任意嵌套异常和用户态栈切换不在上述评分覆盖中，不能凭100分推断全部正确。我们将实际启动和GDB记录留在本地，GitHub只保留代码与报告。


\section{原理联系、局限和总结}

本实验中的现场保存与恢复对应操作系统的上下文管理，但尚未切换进程；时钟提供调度的基础事件，但本次没有进程调度器；非法指令处理体现同步异常分发，但按需分页、页面置换和用户系统调用尚未实现。关闭本核中断保护临界区不等于多核锁，volatile使ticks按要求进行可见访问，也不等于原子读改写。


通过本次实验，我们接通了从汇编入口、C处理函数到sret返回的完整链路，理解了硬件记录控制状态与软件保存通用寄存器的分工。


参考：\href{http://oslab.mobisys.top/lab2026/_book/lab3/lab3_2_1_exercise.html}{课程练习}、\href{http://oslab.mobisys.top/lab2026/_book/lab3/lab3_3_2_trapentry.html}{Trap入口}、\href{http://oslab.mobisys.top/lab2026/_book/lab3/lab3_3_4_clock.html}{时钟}、\href{http://oslab.mobisys.top/lab2026/_book/lab3/lab3_4_lab_requirement.html}{报告要求}。


\end{document}
